\documentclass{article}
\usepackage{graphicx} % Required for inserting images
\usepackage[italian]{babel}
\usepackage{amsmath}
\usepackage[hidelinks]{hyperref}
\usepackage{amsfonts}
\usepackage{amssymb}
\usepackage{accents}
\usepackage{xcolor}
\usepackage{amsthm}
\usepackage{diagbox}


\newtheorem{theorem}{Teorema}
\newtheorem{lemma}{Lemma}

\newcommand{\df}{\noindent\textbf{Definizione} }

\newcommand{\ubar}[1]{\underaccent{\bar}{#1}}


\title{Statistica}
\author{Leonardo Ganzaroli}
\date{}

\begin{document}

\maketitle

\tableofcontents

\newpage

\hypersetup{allcolors=black}


\section{Elementi utili}

\noindent\textbf{Bernoulli} \[ \boxed{X\sim \operatorname{Ber}(p)} \]

\[
\begin{aligned}
P(X=x)
&=
p^x(1-p)^{1-x},
\qquad x\in\{0,1\},
\\[0.5em]
E[X]
&=
p,
\qquad
\operatorname{Var}(X)=p(1-p).
\end{aligned}
\]

\medskip

\noindent\textbf{Binomiale} \[ \boxed{X\sim \operatorname{Bin}(n,p)} \]

\[
\begin{aligned}
P(X=x)
&=
\binom{n}{x}p^x(1-p)^{n-x},
\qquad x=0,\ldots,n,
\\[0.5em]
E[X]
&=
np,
\qquad
\operatorname{Var}(X)=np(1-p).
\end{aligned}
\]

\medskip

\noindent\textbf{Multinomiale} \[ \boxed{ (X_1,\ldots,X_k) \sim \operatorname{Mult}(n,p_1,\ldots,p_k) } \]

\[
\begin{aligned}
P(X_1=x_1,\ldots,X_k&=x_k)
=
\frac{n!}{x_1!\cdots x_k!}
\prod_{i=1}^{k}p_i^{x_i},
\qquad
\sum_{i=1}^{k}x_i=n,
\\[0.5em]
E[X_i]
&=
np_i,\qquad
\operatorname{Var}(X_i)
=
np_i(1-p_i).
\end{aligned}
\]

\medskip

\noindent\textbf{Poisson} \[ \boxed{X\sim\operatorname{Pois}(\lambda)} \]

\[
\begin{aligned}
P(X=x)
&=
\frac{\lambda^x e^{-\lambda}}{x!},
\qquad x=0,1,2,\ldots,\ \lambda>0
\\[0.5em]
E[X]
&=
\operatorname{Var}(X)=\lambda.
\end{aligned}
\]

\medskip

%------------------------------------------------

\noindent\textbf{Uniforme continua} \[ \boxed{X\sim U(a,b)} \]

\[
\begin{aligned}
f(x)
&=
\frac1{b-a},
\qquad a<x<b,
\\[0.5em]
E[X]
&=
\frac{a+b}{2},
\qquad
\operatorname{Var}(X)=\frac{(b-a)^2}{12}.
\end{aligned}
\]

\newpage

\noindent\textbf{Esponenziale} \[ \boxed{X\sim\operatorname{Exp}(\lambda)} \]

\[
\begin{aligned}
f(x)
&=
\lambda e^{-\lambda x},
\qquad x,\ \lambda>0,
\\[0.5em]
E[X]
&=
\frac1\lambda,
\qquad
\operatorname{Var}(X)=\frac1{\lambda^2}.
\end{aligned}
\]

\medskip

\noindent\textbf{Gamma} ($\beta$ parametro di tasso) \[ \boxed{X\sim\Gamma(\alpha,\beta)} \]


\[
\begin{aligned}
f(x)
&=
\frac{\beta^\alpha}{\Gamma(\alpha)}
x^{\alpha-1}e^{-\beta x},
\qquad x>0,\ \alpha,\beta >0,
\\[0.5em]
E[X]
&=
\frac{\alpha}{\beta},
\qquad
\operatorname{Var}(X)=\frac{\alpha}{\beta^2}.
\end{aligned}
\]

Si ha

\[
X_1,\ldots,X_n
\stackrel{\mathrm{iid}}{\sim}
\operatorname{Exp}(\lambda)
\Rightarrow
\sum_{i=1}^{n}X_i
\sim
\Gamma(n,\lambda)
\]

Inoltre

\[
\bar X=\frac1n\sum_{i=1}^{n}X_i
\sim
\Gamma(n,n\lambda)
\]

%------------------------------------------------

\noindent\textbf{Beta} \[ \boxed{X\sim\operatorname{Beta}(\alpha,\beta)} \]

\[
\begin{aligned}
f(x)
&=
\frac{\Gamma(\alpha+\beta)}
{\Gamma(\alpha)\Gamma(\beta)}
x^{\alpha-1}(1-x)^{\beta-1},
\qquad 0<x<1,\ \alpha,\beta >0,
\\[0.5em]
E[X]
&=
\frac{\alpha}{\alpha+\beta},
\qquad
\operatorname{Var}(X)=
\frac{\alpha\beta}
{(\alpha+\beta)^2(\alpha+\beta+1)}.
\end{aligned}
\]

%------------------------------------------------

\noindent\textbf{Normale} \[ \boxed{X\sim N(\mu,\sigma^2)} \]

\[
\begin{aligned}
f(x)
&=
\frac1{\sigma\sqrt{2\pi}}
\exp\left(
-\frac{(x-\mu)^2}{2\sigma^2}
\right),\qquad \sigma>0,
\\[0.5em]
E[X]
&=
\mu,
\qquad
\operatorname{Var}(X)=\sigma^2.
\end{aligned}
\]

\vspace{10pt}

\[
\boxed{
\textbf{N.B.}\qquad
X\sim N(\mu,\sigma^2)
\Rightarrow
Z=\frac{X-\mu}{\sigma}\sim N(0,1)
}
\]

\newpage

Se

\[
X_1,\ldots,X_n \text{ indipendenti},
\qquad
X_i\sim N(\mu_i,\sigma_i^2)
\]

allora

\[
\sum_{i=1}^{n}X_i
\sim
N\left(
\sum_{i=1}^{n}\mu_i,
\sum_{i=1}^{n}\sigma_i^2
\right)
\]

Inoltre

\[
X_1,\ldots,X_n
\stackrel{\mathrm{iid}}{\sim}
N(\mu,\sigma^2)
\Rightarrow
\bar X
\sim
N\left(
\mu,\frac{\sigma^2}{n}
\right)
\]

%------------------------------------------------

\noindent\textbf{t di Student} \[ \boxed{X\sim t_\nu} \]

\[
f(x)
=
\frac{\Gamma\!\left(\frac{\nu+1}{2}\right)}
{\sqrt{\nu\pi}\Gamma\!\left(\frac{\nu}{2}\right)}
\left(1+\frac{x^2}{\nu}\right)^{-\frac{\nu+1}{2}},
\]

\[
E[X]=0
\quad(\nu>1),
\qquad
\operatorname{Var}(X)
=
\frac{\nu}{\nu-2}
\quad(\nu>2).
\]

Se

\[
Z\sim N(0,1),
\qquad
V\sim\chi^2_\nu
\]

sono indipendenti, allora

\[
\frac{Z}{\sqrt{V/\nu}}
\sim t_\nu
\]

Inoltre:

\[
X_1,\ldots,X_n
\stackrel{\mathrm{iid}}{\sim}
N(\mu,\sigma^2)
\Rightarrow
\frac{\bar X-\mu}{S/\sqrt n}
\sim t_{n-1}
\]

%------------------------------------------------

\noindent\textbf{Chi-quadro} \[ \boxed{X\sim\chi^2_k} \]

\[
Z_1,\ldots,Z_k
\stackrel{\mathrm{iid}}{\sim}
N(0,1)
\Rightarrow
\sum_{i=1}^{k}Z_i^2
\sim\chi^2_k,
\]

\[
E[X]=k,
\qquad
\operatorname{Var}(X)=2k.
\]

Si ha

\[
X_1,\ldots,X_n
\stackrel{\mathrm{iid}}{\sim}
N(\mu,\sigma^2)\Rightarrow
\boxed{
\frac{(n-1)S^2}{\sigma^2}
\sim
\chi^2_{n-1}
}
\]\vspace{10pt}

\noindent\textit{Nota: $S^2$ è la varianza campionaria, vedere la sezione 2.}

\newpage

%------------------------------------------------

\noindent\textbf{F di Fisher} \[ \boxed{X\sim F_{d_1,d_2}} \]

Se

\[
A\sim\chi^2_{d_1},
\qquad
B\sim\chi^2_{d_2}
\]

sono indipendenti, allora

\[
X=
\frac{A/d_1}{B/d_2}
\sim
F_{d_1,d_2}
\]

\[
E[X]
=
\frac{d_2}{d_2-2},
\qquad d_2>2,\qquad
\operatorname{Var}(X)
=
\frac{2d_2^2(d_1+d_2-2)}
{d_1(d_2-2)^2(d_2-4)},
\qquad d_2>4.
\]

%------------------------------------------------

\noindent\textbf{Legge (debole) dei Grandi Numeri}:\newline

Se

\[
X_1,\ldots,X_n
\stackrel{\mathrm{iid}}{\sim}X
\]

con

\[
E[|X|]<\infty,
\qquad
E[X]=\mu
\]

allora

\[
\boxed{
\bar X
=
\frac1n\sum_{i=1}^{n}X_i
\overset{P}{\longrightarrow}
\mu
}
\]

\vspace{5pt}

%------------------------------------------------

\noindent\textbf{Teorema del Limite Centrale}:\newline

Se

\[
X_1,\ldots,X_n
\stackrel{\mathrm{iid}}{\sim}X
\]

con

\[
E[X]=\mu,
\qquad
\operatorname{Var}(X)=\sigma^2<\infty
\]

allora:

\[
\boxed{
\frac{\bar X-\mu}{\sigma/\sqrt n}
\overset{d}{\longrightarrow}
N(0,1)
}
\]

ovvero con $n$ grande:

\[
\boxed{
\bar X
\approx
N\left(
\mu,\frac{\sigma^2}{n}
\right)
}
\]



\section{Campioni e statistiche}

\df Un campione casuale di taglia $n$ è una sequenza di variabili aleatorie iid con distribuzione $F_\theta$ e $\theta$ parametro ignoto:

$$X_1,X_2,\cdots, X_n\stackrel{\mathrm{iid}}{\sim}F_\theta$$\newline

Alcuni metodi per generare campioni casuali:
\begin{itemize}
    \item \textbf{Metodo dell'inversione}

    Data la funzione di ripartizione $F(x)=P(X\leq x)$;
    
    Si genera un numero casuale $u$ da $U(0,1)$;
    
    Si calcola $x=F^{-1}(u)$;

    Il numero $x$ sarà distribuito secondo $F$.

    \item \textbf{Trasformazione di Box-Muller (originale)}

    Prendo $U_1,U_2\sim U(0,1)$;

    Con le seguenti formule si ottengono due gaussiane standard indipendenti:
    \[
    Z_1=\sqrt{-2\ln{(U_1)}}\cos(2\pi U_2),\qquad
    Z_2=\sqrt{-2\ln{(U_1)}}\sin(2\pi U_2)
    \]
    
    Con cui si può costruire una normale con media e varianza scelti:
    
    \[
    X=\mu+\sigma Z_{(1 \text{ oppure } 2)}\Rightarrow X\sim N(\mu,\sigma^2)
    \]

    \item \textbf{Metodo di accettazione-rifiuto}

    Data $f(x)$, si sceglie $g(x)$ tale che $\forall x\ f(x)\leq M g(x)\ \ (M\geq\sup_x\frac{f(x)}{g(x)})$;
    
    Si generano $x\sim g(x)$ e $ u\sim U(0,1)$;
    
    Si accetta il valore se $u\leq \frac{f(x)}{M g(x)}$.\newline

\end{itemize}

\df Una statistica è una funzione del campione $(X_1,\cdots,X_n)$:

$$T:\mathbb{R}^n\rightarrow\mathbb{R}$$

Dato che il campione è formato da v.a. la statistica in sé è una v.a.\newline

Alcune statistiche sono:
\begin{itemize}
    \item \textbf{Media campionaria}
    \[
        \bar X=\frac1n\sum_{i=1}^n X_i
    \]

    \item \textbf{Varianza campionaria}
    \[
        S^2=\frac1{n-1}\sum_{i=1}^n (X_i-\bar X)^2
    \]

    \item \textbf{Statistiche d'ordine}

    Ottenute ordinando il campione in ordine crescente:
    \[
        X_{(1)}\le X_{(2)}\le\cdots\le X_{(n)}
    \]
    dove \(X_{(i)}\) rappresenta la \(i\)-esima osservazione più piccola del campione.

    \item \textbf{Mediana campionaria}
    \[
        M=
        \begin{cases}
            X_{\left(\frac{n+1}{2}\right)} & \text{se } n \text{ è dispari},\\[1ex]
            \dfrac{X_{\left(\frac n2\right)}+X_{\left(\frac n2+1\right)}}{2}
            & \text{ altrimenti}.
        \end{cases}
    \]

    \item \textbf{Minimo} e \textbf{Massimo}
    \[
        X_{(1)},\ X_{(n)}
    \]

\end{itemize}

\vspace{7pt}

Alcune proprietà che può avere una statistica:
\begin{description}
    \item[Sufficienza] 

        La statistica contiene tutta l'informazione presente nel campione riguardo a $\theta$.
    
    \item[Minimalità] 

        La statistica è funzione di qualsiasi altra statistica sufficiente.
    
    \item[Completezza]

        Non esistono funzioni non banali della statistica che hanno valore atteso 0 per ogni valore assumibile da $\theta$.

    \item[Ancillarità]

        La distribuzione della statistica non dipende dal parametro da stimare.
    
\end{description}

\vspace{5pt}

\begin{theorem}[Basu]
    Qualsiasi statistica sufficiente e limitatamente completa\footnote{La completezza limitata è una forma più debole della completezza.} è indipendente da qualsiasi statistica ancillare.
\end{theorem}

\vspace{5pt}

Esistono dei metodi per capire se una statistica ha le prime 2 proprietà:
\begin{itemize}
    \item \textbf{Sufficienza}
        \begin{itemize}
            \item Capire se $P_\theta(X=x\ |\ T(X)=t)$ non dipende da $\theta$;
            \item Usare la fattorizzazione di Neyman-Fisher:
                $$\prod_{i=1}^nf_\theta(x_i)=g(T(x),\theta)h(x)$$

                Ossia capire se si può ``dividere'' la produttoria in una parte che dipende dal parametro e dalla statistica e un'altra che dipende solo dal campione.
                
        \end{itemize}
    \item \textbf{Minimalità}
            $$\text{Verificare che (dati 2 campioni $x,y$): }
            \frac{\prod_{i=1}^nf_\theta(x_i)}{\prod_{i=1}^nf_\theta(y_i)} \text{ non dipende da } \theta \iff T(X)=T(Y)$$
\end{itemize}

Esempi:

\begin{itemize}
    \item \textbf{Esponenziale con $\lambda$ ignota} 

        $$\lambda e^{-\lambda x}\rightarrow\prod_{i=1}^n\lambda e^{-\lambda x_i}\rightarrow\textcolor{red}{1}*\textcolor{blue}{\lambda^n e^{-\lambda \sum_{i=1}^nx_i}}$$

        In questo caso l'intera formula è \textcolor{blue}{$g(T(X),\lambda)$}, moltiplicando la formula per 1 posso prendere quest'ultimo come \textcolor{red}{$h(x)$}. Quindi la somma è sufficiente.

    \item \textbf{Uniforme $(0,\theta)$} 

        (Qui devo considerare anche il vincolo dato che in questo compare $\theta$)

        $$\frac{1}{\theta},\ 0<x<\theta\rightarrow\prod_{i=1}^n\frac{1}{\theta}\mathbf{1}_{(0<x_i<\theta)}\rightarrow\frac{1}{\theta^n}\mathbf{1}_{(x_{(n)}<\theta)}$$

        (Il prodotto delle funzioni indicatrici vale 1 solamente se ogni $x_i$ è minore di $\theta$, in modo equivalente vale 1 solo se $\theta$ è maggiore dell'elemento massimo del campione)

        Mettendo a rapporto:

        \[
            \frac{\frac{1}{\theta^n}\mathbf{1}_{(x_{(n)}<\theta)}}
            {\frac{1}{\theta^n}\mathbf{1}_{(y_{(n)}<\theta)}}
            =
            \frac{\mathbf{1}_{(x_{(n)}<\theta)}}
            {\mathbf{1}_{(y_{(n)}<\theta)}} 
        \]

        Qui l'unico caso in cui non compare $\theta$ è se la frazione equivale ad 1, cosa che succede solo se il massimo dei due campioni è lo stesso valore. Quindi $x_{(n)}$ è sufficiente minimale.
    
\end{itemize}

\newpage

\section{Stimatori}

\df Uno stimatore (puntuale) è una statistica che, a partire dai valori del campione, fornisce una stima di un parametro ignoto.

\subsection{MLE}

Lo stimatore di massima verosimiglianza si ottiene con la seguente formula:

$$\hat{\theta}_{MLE}=\arg\max_{\theta\in\Theta}\prod_{i=1}^{n} f_{\theta}(x_i)$$

Spesso conviene lavorare con la versione logaritmica:

$$\arg\max_{\theta\in\Theta}\sum_{i=1}^{n}\log f_{\theta}(x_i)$$

\vspace{10pt}

Un metodo comune per trovare il massimo è usare le derivate:

$$\frac{x^{\frac{1-\theta}{\theta}}}{\theta}\rightarrow\theta^{-n}(\prod_{i=1}^{n}x_i)^{\frac{1-\theta}{\theta}}\overset{\text{log}}{\longrightarrow}-n\log{\theta}+\theta^{-1}\sum\log{x_i}-\sum\log{x_i}$$

\begin{align*}
    \frac{dl}{d\theta} &\rightarrow-\frac{n}{\theta}-\frac{\sum\log{x_i}}{\theta^2}\overset{=0}{\longrightarrow}\hat{\theta}=-\frac{\sum\log{x_i}}{n}\\
    \frac{d^2l}{d\theta^2} &\rightarrow\frac{n}{\theta^2}+\frac{2\sum\log{x_i}}{\theta^3}\\
    &=  \frac{n}{\theta^2}+\frac{2(-n\theta)}{\theta^3}\ \text{(nel punto critico } \sum\log{x_i}=-n\theta)\\
    &= \frac{n}{\theta^2}\left(1+\frac{2(-n\theta)}{n\theta}\right) = -\frac{n}{\theta^2}<0\\
\end{align*}

Quindi il punto critico è un massimo e di conseguenza MLE.\newline

\[
\boxed{
\begin{gathered}
\textbf{N.B.\footnotemark}\quad
\text{Per $n$ grande e sotto opportune condizioni di regolarità:}\\[0.5em]
\hat{\theta}_{\mathrm{MLE}}
\overset{a}{\sim}
N\!\left(\theta_{\text{vero}},\frac{1}{nI(\theta)}\right)
\end{gathered}
}
\]

\footnotetext{La definizione di $I(\theta)$ è nella sezione 3.4.}

\newpage

\subsection{Metodo dei momenti}

Dati $k$ parametri ignoti, gli stimatori dei momenti si ottengono uguagliando i momenti teorici\footnote{I momenti teorici sono riferiti alla singola variabile aleatoria $X$, non al campione.} della distribuzione con i relativi momenti campionari:

\begin{align*}
E[X] &= m_1\\
E[X^2] &= m_2\\
&\vdots\\
E[X^k] &= m_k\\
\\
\text{Dove }m_r=\frac1n\sum_{i=1}^n X_i^r,&\qquad r=1,\dots,k.
\end{align*}

Per esempio una normale con entrambi i parametri ignoti:

\begin{align*}
    E[X]=\mu,&\qquad E[X^2]=\mu^2+\sigma^2\\
    \bar X=\frac1n\sum_{i=1}^n X_i,&\qquad \bar X^2=\frac1n\sum_{i=1}^n X_i^2\\
    \\
    E[X]=m_1
    &\longrightarrow
    \hat\mu_{MM}=\bar X,\\
    E[X^2]=m_2
    &\longrightarrow
    \mu^2+\sigma^2=\frac1n\sum_{i=1}^nX_i^2\\
    &\longrightarrow
    \hat\sigma^2_{MM}
    =\frac1n\sum_{i=1}^nX_i^2-\bar X^2\ (\text{Uso $\hat\mu_{MM}$})
\end{align*}

\[
\boxed{
\parbox{0.85\textwidth}{
\textbf{N.B.} Il metodo dei momenti può presentare problemi:

\begin{itemize}\setlength{\itemsep}{0pt}
    \item Per la distribuzione di Cauchy il primo momento non esiste, quindi
    questo metodo non è applicabile;
    
    \item Per 
    \[
    \frac{1+\theta\cos x}{2\pi}
    \]
    si ha $E[X]=\pi$ (non compare $\theta$), bisogna quindi utilizzare un momento
    di ordine superiore.
\end{itemize}
}
}
\]

\subsection{Bayes}

\df La distribuzione a priori ($\pi(\theta)$) è una distribuzione di probabilità assegnata al parametro prima di osservare i dati\footnote{Quando una distribuzione viene utilizzata come prior, la variabile generica della densità (comunemente $x$) viene indicata con il simbolo del parametro stesso.}.\newline 

\df La verosimiglianza è la probabilità (o densità) del campione osservato vista come funzione del parametro: \[ L(\theta;x)=\prod_{i=1}^{n}f_\theta(x_i). \]

\df La distribuzione a posteriori (posterior) è la distribuzione del parametro dopo aver osservato il campione, è data da:

\[
\pi(\theta\mid x)\propto
L(\theta;x)\pi(\theta)
\]\newline

Per ottenere lo stimatore bayesiano si sceglie una funzione di perdita\footnote{Qui $\theta$ rappresenta il vero valore del parametro, $a$ la stima.}:

\begin{itemize}
\item \textbf{$(\theta-a)^2$} (MSE):

    Lo stimatore è la media della distribuzione a posteriori;

\item \textbf{$|\theta-\hat{\theta}|$}:

    Lo stimatore è la mediana della distribuzione a posteriori;

\item \textbf{$\arg\max_\theta \pi(\theta\mid x)$}\footnote{Per parametri discreti.}:

    Lo stimatore è il valore che massimizza la posterior;

\end{itemize}

\vspace{6pt}

Esempio:

\begin{align*}
&f(x\mid\theta)=\theta^{x}(1-\theta)^{1-x},
\quad 0<\theta<1,\qquad
\pi(\theta)=\frac{\Gamma(\alpha+\beta)}
{\Gamma(\alpha)\Gamma(\beta)}
\theta^{\alpha-1}(1-\theta)^{\beta-1}\\
\\
&L(\theta;x)
=\prod_{i=1}^{n}
\theta^{x_i}(1-\theta)^{1-x_i}=
\theta^{\sum_{i=1}^{n}x_i}
(1-\theta)^{\sum_{i=1}^{n}(1-x_i)}=
\theta^{\sum_{i=1}^{n}x_i}
(1-\theta)^{n-\sum_{i=1}^{n}x_i}\\
\\
&\pi(\theta\mid x)
\propto
\theta^{\sum x_i}
(1-\theta)^{n-\sum x_i}
\theta^{\alpha-1}
(1-\theta)^{\beta-1}
\propto
\theta^{\alpha+\sum x_i-1}
(1-\theta)^{\beta+n-\sum x_i-1}\\
\\
&\text{Che è una }\mathrm{Beta}
\left(
\alpha+\sum_{i=1}^{n}x_i,\,
\beta+n-\sum_{i=1}^{n}x_i
\right), \text{ scegliendo } MSE\rightarrow \hat{\theta}_B=\frac{\alpha+\sum_{i=1}^{n}x_i}{\alpha+\beta+n}
\end{align*}\newline

\subsection{Proprietà e teoremi vari}

\df L'informazione di Fisher\footnote{Si riferisce alla singola variabile aleatoria $X$, per avere quella del campione basta moltiplicare per la taglia dello stesso.} è definita come:

$$I(\theta)=E_\theta\left[\left(\frac{\partial}{\partial\theta}\log f_\theta(x)\right)^2\right]$$

Oppure\footnote{Sotto opportune condizioni di regolarità e se la funzione è doppiamente differenziabile.}:

$$-E_\theta\left[\frac{\partial^2}{\partial\theta^2}\log f_\theta(x)\right]$$\newline

\df Lo squarto quadratico medio è definito come:

$$MSE(\hat{\theta})=Var_\theta(\hat{\theta})+(E_\theta[\hat{\theta}]-\theta)^2$$\newline

\df La disuguaglianza di Cramér–Rao è definita come:

$$Var(\hat{\theta})\geq\frac{\left(\frac{d}{d\theta}E[\hat{\theta}]\right)^2}{nI(\theta)}$$

\vspace{10pt}

Alcune proprietà che può avere uno stimatore:
\begin{itemize}
    \item \textbf{Correttezza}\footnote{Uno stimatore corretto viene anche chiamato non distorto.}

        $$E[\hat{\theta}]=\theta$$

        (Per uno stimatore corretto il secondo termine di $MSE$ non compare e la parte dx di Cramér–Rao diventa $1/nI(\theta)$)
    
    \item \textbf{Consistenza}

        $$\hat{\theta}\overset{P}{\rightarrow}\theta \text{ quando } n \rightarrow\infty$$
    
    \item \textbf{Efficienza}

        Raggiunge il limite inferiore di Cramér–Rao ed è corretto.\newline
    
\end{itemize}

\df Uno stimatore è detto \textbf{UMVU} se non è distorto ed ha varianza minima tra tutti gli stimatori non distorti, per ogni dal valore del parametro.

\[
\boxed{
\textbf{N.B.}\qquad \text{Stimatore efficiente } \Rightarrow \text{ Stimatore UMVU}
}
\]

\vspace{10pt}

\begin{theorem}[Lehmann–Scheffé]
    Se uno stimatore non è distorto ed è funzione di una statistica sufficiente e completa, allora è lo stimatore UMVU del parametro.
\end{theorem}

\begin{theorem}[Rao–Blackwell]
    Dato uno stimatore $\hat{\theta}$ ed una statistica sufficiente $T$:

    \[
        \hat{\theta}_{RB}=E[\hat{\theta}\mid T]
    \]

    è uno stimatore per cui:

    $$Var(\hat{\theta}_{RB})\leq Var(\hat{\theta}),\quad E[\hat{\theta}_{RB}]=E[\hat{\theta}]$$

    \vspace{5pt}

    Per il teorema di Lehmann–Scheffé si ha inoltre che se $\hat{\theta}$ è corretto e $T$ è anche completa $\ \hat{\theta}_{RB}$ è UMVU.
    
\end{theorem}

\begin{proof}
    $\ $\newline
    $\ $\newline
    Per la legge delle aspettative iterate:
    $$E[\hat{\theta}_{RB}]=E[E[\hat{\theta}\mid T]]=E[\hat{\theta}]$$
    Per la legge della varianza totale:
    $$\textcolor{olive}{Var(\hat{\theta})}=E[Var(\hat{\theta}\mid T)]+Var(E[\hat{\theta}\mid T])=\textcolor{red}{E[Var(\hat{\theta}\mid T)]}+\textcolor{blue}{Var(\hat{\theta}_{RB})}$$
    Dove \textcolor{red}{il primo termine} è maggiore o uguale a 0, quindi \textcolor{blue}{il secondo} è minore o uguale alla \textcolor{olive}{varianza di $\hat{\theta}$}.
\end{proof}

\newpage

\section{Test d'ipotesi e intervalli di confidenza}

\subsection{IC}

\df Una quantità pivotale $(Q(X,\theta))$ è una funzione del campione e del parametro la cui distribuzione non dipende dal valore del parametro.\newline

\df Un intervallo di confidenza per un parametro è un intervallo costruito partendo dal campione, tale che:

\[
P(u(X)<\theta<v(X))=\gamma
\]

\vspace{5pt}

\noindent Dove $\gamma$ è detto livello di confidenza. Con un livello di confidenza del 95\%, se si ripete molte volte il campionamento e per ogni campione si costruisce un intervallo con la stessa procedura, circa il 95\% degli intervalli ottenuti conterrà il vero valore del parametro.\newline

\df Il livello di significatività è definito come $\alpha=1 - \gamma$.\newline

Un metodo comune per costruire questi intervalli consiste nello sfruttare le quantità pivotali, ad esempio per una normale:

$$\frac{\overline{X}-\mu}{\sigma/\sqrt{n}}\sim N(0,1),\quad \frac{\overline{X}-\mu}{S/\sqrt{n}}\sim t_{n-1},\quad \frac{(n-1)S^2}{\sigma^2}\sim\chi^2_{n-1}$$

Si possono usare rispettivamente per costruire gli intervalli di:
\begin{itemize}

    \item \textbf{Media (varianza nota):}
    
    \[
    \frac{\overline{X}-\mu}{\sigma/\sqrt{n}}\sim N(0,1)
    \quad\Longrightarrow\quad
    P\left(-z_{1-\frac{\alpha}{2}}\leq
    \frac{\overline{X}-\mu}{\sigma/\sqrt{n}}
    \leq z_{1-\frac{\alpha}{2}}\right)=1-\alpha
    \]
    
    \[
    \Longrightarrow\quad
    \mu\in
    \left(
    \overline{X}-z_{1-\frac{\alpha}{2}}\frac{\sigma}{\sqrt n},
    \overline{X}+z_{1-\frac{\alpha}{2}}\frac{\sigma}{\sqrt n}
    \right)
    \]

    \item \textbf{Media (varianza ignota):}
    
    \[
    \frac{\overline{X}-\mu}{S/\sqrt{n}}\sim t_{n-1}
    \quad\Longrightarrow\quad
    P\left(-t_{1-\frac{\alpha}{2},n-1}\leq
    \frac{\overline{X}-\mu}{S/\sqrt{n}}
    \leq t_{1-\frac{\alpha}{2},n-1}\right)=1-\alpha
    \]
    
    \[
    \Longrightarrow\quad
    \mu\in
    \left(
    \overline{X}-t_{1-\frac{\alpha}{2},n-1}\frac{S}{\sqrt n},
    \overline{X}+t_{1-\frac{\alpha}{2},n-1}\frac{S}{\sqrt n}
    \right)
    \]

    \item \textbf{Varianza:}
    
    \[
    \frac{(n-1)S^2}{\sigma^2}\sim\chi^2_{n-1}
    \quad\Longrightarrow\quad
    P\left(
    \chi^2_{\frac{\alpha}{2},n-1}
    \leq
    \frac{(n-1)S^2}{\sigma^2}
    \leq
    \chi^2_{1-\frac{\alpha}{2},n-1}
    \right)=1-\alpha
    \]
    
    \[
    \Longrightarrow\quad
    \sigma^2\in
    \left(
    \frac{(n-1)S^2}{\chi^2_{1-\frac{\alpha}{2},n-1}},
    \frac{(n-1)S^2}{\chi^2_{\frac{\alpha}{2},n-1}}
    \right)
    \]\newline

\end{itemize}

Esempi:
\begin{itemize}

    \item \textbf{Media:}

    \[
        \overline{X}=12,\qquad n=100,\qquad S=4,\qquad \alpha=0.05,\qquad \sigma=4
    \]

    \begin{itemize}

    \item \textbf{Varianza nota:}

    \[
    z_{1-\frac{\alpha}{2}}=z_{0.975}=1.96\rightarrow
    12\pm1.96\frac{4}{\sqrt{100}}
    =
    12\pm0.784
    \]

    \[
    IC_{95\%}=(11.216,\;12.784)
    \]


    \item \textbf{Varianza ignota:}

    Ipotizziamo $S=4$

    \[
    t_{1-\frac{\alpha}{2},n-1}
    =
    t_{0.975,99}
    \approx1.984\rightarrow
    12\pm1.984\frac{4}{\sqrt{100}}
    =
    12\pm0.794
    \]

    \[
    IC_{95\%}=(11.206,\;12.794)
    \]

    \end{itemize}


    \item \textbf{Varianza:}

    \[
    S^2=9,\qquad n=20,\qquad \alpha=0.05
    \]

    \[
    \chi^2_{0.975,19}=32.852,
    \quad
    \chi^2_{0.025,19}=8.907\rightarrow
    \left(
    \frac{19\cdot9}{32.852},
    \frac{19\cdot9}{8.907}
    \right)
    \]

    \[
    IC_{95\%}=(5.20,\;19.20)
    \]

\end{itemize}

\subsection{Test}

\df Un test d'ipotesi è una procedura che permette di decidere, basandosi sul campione, se rifiutare una certa ipotesi sul parametro. Si hanno 2 ipotesi:
\begin{itemize}
    \item $H_0$ (ipotesi nulla, si considera vera fino a prova contraria)
    
    \item $H_1$ (ipotesi alternativa)
    
\end{itemize}


\[
    \Theta=\Theta_0 \cup \Theta_1,\quad \Theta_0=\{\theta\in\Theta\mid H_0\text{ è vera}\},\quad \Theta_1=\{\theta\in\Theta\mid H_1\text{ è vera}\}
\]

\vspace{8pt}

\df La regione critica ($C$) è l'insieme di valori della statistica di test per cui $H_0$ viene rifiutata.\newline

\df La funzione di potenza è la probabilità di rifiutare $H_0$ al variare del parametro.

\[
P_\theta(T(X)\in C)
\]\newline

\df L'errore di I specie consiste nel rifiutare $H_0$ quando essa è vera\footnote{$\alpha$ è il livello di significatività.}:

\[
\sup_{\theta\in\Theta_0}P_\theta(T(X)\in C)\leq\alpha
\]\newline

\df L'errore di II specie consiste nel non rifiutare $H_0$ quando essa è falsa:

\[
\beta
=
P_\theta(T(X)\notin C)
\qquad
\theta\in\Theta_1.
\]\newline

\df La potenza del test è $1-\beta$.\newline

\[
\boxed{\alpha\downarrow\ \Rightarrow\beta\uparrow\ \Rightarrow\text{ Potenza test}\downarrow}
\]\newline

\df Un test di livello $\alpha$ è \textbf{MP} se, tra tutti i test dello stesso livello, ha la funzione di potenza massima per un valore fissato $\theta_1\in\Theta_1$.\newline

\df Un test di livello $\alpha$ è \textbf{UMP} se, tra tutti i test dello stesso livello, ha la funzione di potenza massima per ogni $\theta\in\Theta_1$.\newpage

Tipi di test:
\begin{itemize}
    \item \textbf{Test bilaterale:}
    \[
    H_0:\theta=\theta_0,
    \qquad
    H_1:\theta\neq\theta_0
    \]
    La regione critica è divisa in due code della distribuzione.

    \item \textbf{Test unilaterale destro:}
    \[
    H_0:\theta\leq\theta_0,
    \qquad
    H_1:\theta>\theta_0
    \]
    La regione critica si trova nella coda destra della distribuzione.

    \item \textbf{Test unilaterale sinistro:}
    \[
    H_0:\theta\geq\theta_0,
    \qquad
    H_1:\theta<\theta_0
    \]
    La regione critica si trova nella coda sinistra della distribuzione.\newline
\end{itemize}

\begin{lemma}[Neyman-Pearson]
Per  due ipotesi semplici $H_0:\theta=\theta_0,\ H_1:\theta=\theta_1$ il test che rifiuta $(H_0)$ quando:

\[
\Lambda(X)=\frac{L(\theta_0;X)}{L(\theta_1;X)}\leq k\quad
\text{dove $k$ è  tale che }
P_{\theta_0}(\Lambda(X)\leq k)=\alpha
\]

è MP.\newline
\end{lemma}

Esempi:
\begin{itemize}
    \item Un'azienda dichiara che il peso medio delle confezioni è di $500$ g, si sospetta però che il peso medio sia maggiore.

    
Da un campione di $n=36$ confezioni si ottiene $
\bar X=504,\ \sigma=12$. Si vuole un test con $\alpha=0.05$

\[
H_0:\mu\leq500,
\qquad
H_1:\mu>500.
\]

Essendo la varianza nota, si può usare:

\begin{align*}
Z&=\frac{\bar X-\mu_0}{\sigma/\sqrt n}=\frac{504-500}{12/\sqrt{36}}
=
\frac{4}{2}
=
2 \\
\\
C&=\{Z>z_{1-\alpha}\}=\{Z>1.645\}=\{2>1.645\}\rightarrow\text{ Rifiuto }H_0
\end{align*}

\item Si vuole verificare se il tempo medio di consegna di un servizio sia diverso
da $30$ minuti.

Da un campione di $n=16$ osservazioni si ottiene $\bar X=32,\
S=5$. Si vuole un test con $\alpha=0.05$

\[
H_0:\mu=30,
\qquad
H_1:\mu\neq30.
\]

La varianza è ignota, quindi si usa:

\begin{align*}
T&=\frac{\bar X-\mu_0}{S/\sqrt n}=
\frac{32-30}{5/\sqrt{16}}
=
\frac{2}{1.25}
=
1.6
\\\\
C&=
\{|T|>t_{15,1-\alpha/2}\}=\{|T|>2.131\}=\{|1.6|>2.131\}\rightarrow\text{ Non rifiuto }H_0
\end{align*}

\item Si presume una deviazione standard $\sigma_0=4$.


Da un campione normale di dimensione $n=20$ si ottiene $S^2=25$. Si vuole un test con $\alpha=0.05$

\[
H_0:\sigma^2=16,
\qquad
H_1:\sigma^2\neq16.
\]

\begin{align*}
\chi^2&=
\frac{(n-1)S^2}{\sigma_0^2}=
\frac{19\cdot25}{16}
=
29.69
\\
\\
C&=
\{\chi^2<\chi^2_{19,0.025}\}
\cup
\{\chi^2>\chi^2_{19,0.975}\}=\{\chi^2<8.907\}\cup\{\chi^2>32.852\}\rightarrow\text{ Non rifiuto }H_0
\end{align*}

\item Ho $(1-\theta)^{x-1}\theta$, dato un campione di taglia 1 $(x=4)$ e una significatività del $5\%$:

$$H_0:\theta<1\%\ ?$$

Questa è una distribuzione geometrica, più cresce $\theta$ meno prove dovrò fare per avere un successo. Quindi $C=\{x\leq c\}$.

\begin{align*}
    \sup_{\theta\leq0.01}P_\theta(X\leq c)&=P_{0.01}(X\leq c)\leq0.05\qquad(P_p(X\geq c+1)=(1-p)^c)\\
    &=1-(1-0.01)^c\leq0.05\rightarrow1-0.99^c\leq0.05\\
    &\rightarrow0.99^c\geq0.95\rightarrow c\leq\frac{\ln0.95}{\ln0.99}\rightarrow c\leq5.12\rightarrow c=5
\end{align*}

Quindi $C=\{x\leq5\}$, rifiuto $H_0$.

\item Una moneta viene lanciata $10$ volte. Si vuole verificare se la moneta è equa. (Significatività del $5\%$)

\[
H_0:p=0.5,
\qquad
H_1:p=0.8.
\]

Uso il lemma di Neyman-Pearson.

\begin{align*}
&
\frac{\binom{10}{x}(0.5)^x(0.5)^{10-x}}
{\binom{10}{x}(0.8)^x(0.2)^{10-x}}=
\\
&=
\frac{(0.5)^{10}}
{(0.8)^x(0.2)^{10-x}}=
\\
&=
\left(\frac{0.5}{0.8}\right)^x
\left(\frac{0.5}{0.2}\right)^{10-x}=
\\
&=
(0.625)^x(2.5)^{10-x}.
\end{align*}

\[
P_{0.5}(X\geq k)\leq0.05\ \text{ (Il rapporto decresce rispetto ad $x$)}
\]

\begin{align*}
P_{0.5}(X\geq9)
&=
P(X=9)+P(X=10)
\\
&=
\binom{10}{9}(0.5)^{10}
+
\binom{10}{10}(0.5)^{10}
\\
&=
\frac{10+1}{1024}=
0.0107<0.05
\end{align*}

\[
P_{0.5}(X\geq8)
=
\frac{45+10+1}{1024}
=
0.0547>0.05.
\]

Quindi $C=\{X\geq9\}$.

\end{itemize}

\[
\fbox{%
\begin{minipage}{0.85\textwidth}
Si noti come esiste una correlazione con gli intervalli di confidenza. Nel terzo esempio, costruendo un intervallo di confidenza con livello
\(1-\alpha = 95\%\), si ottiene
\[
(14.46,\;52.78)
\]
che contiene il valore \(16\), ossia il valore ipotizzato sotto \(H_0\).

Questo non è un caso: gli intervalli di confidenza contengono esattamente i valori del parametro per i quali il corrispondente test d'ipotesi non rifiuta \(H_0\).
\end{minipage}%
}
\]

\vspace{7pt}

\df Il p-value è un valore che misura quanto i dati osservati siano compatibili con l'ipotesi nulla: un valore piccolo indica che il risultato osservato è poco plausibile se $H_0$ fosse vera.


\section{\texorpdfstring{Test $\chi^2$}{Chi-quadro}}

\subsection{Bontà di adattamento}

\df Una variabile categorica è una variabile che assume valori rappresentati da categorie, utilizzate per classificare le osservazioni in gruppi distinti.\newline

Data una variabile con $k$ categorie e un campione di
dimensione $n$, si vuole testare:

\[
H_0:
(O_1,\ldots,O_k)\sim
\operatorname{Mult}(n,p_1,\ldots,p_k)\text{ (con $O_i$ osservazioni della categoria $i$)}
\]\newline

La statistica usata è: $$
\sum_{i=1}^{k}
\frac{(O_i-E[X_i])^2}{E[X_i]}\overset{\text{a}}{\sim}\chi^2_{k-m-1} \text{ ($m$ è il numero di parametri da stimare e }E[X_i]=np_i)
$$

La convergenza vale sotto $H_0$, per $n$ grande e se $\forall i\ E[X_i]\geq5$.\newline

\[
\text{La regione critica è:}
\left\{
\text{Valore statistica}>
\chi^2_{1-\alpha,k-m-1}
\right\}
\]\newline

Esempio (5\%):

\begin{table}[ht]
    \centering
    \begin{tabular}{c|c|c}
        $O_1$ & $O_2$ & $O_3$\\
        \hline
        10 & 68 & 112\\
    \end{tabular}
\end{table}

Con ($p_1=(1-\theta)^2,\ p_2=2\theta(1-\theta),\ p_3=\theta^2$):

\vspace{8pt}

\[
L=
\frac{190!}{10!\,68!\,112!}
((1-\theta)^2)^{10}
(2\theta(1-\theta))^{68}
(\theta^2)^{112}
\]

\[
\frac{d l}{d\theta}=0
\quad\longrightarrow\quad
\hat\theta=
\frac{68+2(112)}{2\cdot190}
=0.768
\]

\vspace{8pt}

\[
\begin{aligned}
E_1 &= 190(1-0.768)^2 \approx 10.18,\\
E_2 &= 190\bigl(2\cdot0.768(1-0.768)\bigr) \approx 67.68,\\
E_3 &= 190(0.768)^2 \approx 112.14
\end{aligned}
\]

\vspace{8pt}

$$\frac{(10-10.18)^2}{10.18}+\frac{(68-67.68)^2}{67.68}+\frac{(112-112.14)^2}{112.14}\approx0.005$$

\vspace{7pt}

$$\chi^2_{0.95,1}=3.84>0.005\rightarrow\text{ Non rifiuto $H_0$}$$

\subsection{Indipendenza}

\df Una tabella di contingenza è una tabella a doppia entrata che riporta le frequenze congiunte di due (o più) variabili categoriche, mostrando quante osservazioni appartengono a ciascuna combinazione di categorie.\newline

Date due variabili categoriche, inserite in una tabella con $r$ righe e $c$ colonne. Si vuole testare:

\[
H_0:\text{ le due variabili sono indipendenti}
\]\newline

La statistica del test è:

$$
X^2=
\sum_{i=1}^{r}
\sum_{j=1}^{c}
\frac{(O_{ij}-E_{ij})^2}{E_{ij}}
\overset{\text{a}}{\sim}
\chi^2_{(r-1)(c-1)}
$$$$
\left(\text{con $O_{ij}$ frequenza osservata nella cella $(i,j)$}
\text{ e }E_{ij}=
\frac{(\text{totale riga }i)(\text{totale colonna }j)}{n}\right)
$$\newline

Valgono gli stessi criteri del test precedente per la convergenza.\newline

\[
\text{La regione critica è:}
\left\{
\text{Valore statistica}>
\chi^2_{1-\alpha,(r-1)(c-1)}
\right\}
\]\newline

Esempio (10\%):

\begin{table}[ht]
    \centering
    \begin{tabular}{c|c|c|c|c}
        \diagbox{$A$}{$B$} & $B_1$ & $B_2$ & $B_3$ & Tot.\\
        \hline
        $A_1$ & 15 & 23 & 10 & 48\\
         \hline
        $A_2$ & 25 & 19 & 8 & 52\\
         \hline
        Tot. & 40 & 42 & 18 & 100\\
    \end{tabular}
\end{table}

\[
\begin{aligned}
E_{11}&=\frac{48\cdot40}{100}=19.2, &
E_{12}&=\frac{48\cdot42}{100}=20.16, &
E_{13}&=\frac{48\cdot18}{100}=8.64,\\
E_{21}&=\frac{52\cdot40}{100}=20.8, &
E_{22}&=\frac{52\cdot42}{100}=21.84, &
E_{23}&=\frac{52\cdot18}{100}=9.36.
\end{aligned}
\]

\vspace{8pt}

\[
\begin{aligned}
\frac{(15-19.2)^2}{19.2}
+\frac{(23-20.16)^2}{20.16}
+\frac{(10-8.64)^2}{8.64}+
\frac{(25-20.8)^2}{20.8}
+\frac{(19-21.84)^2}{21.84}
+\frac{(8-9.36)^2}{9.36}
\approx2.66
\end{aligned}
\]

\vspace{8pt}

\[
\chi^2_{0.90,2}
=
4.61
>
2.66\rightarrow\text{ Non rifiuto $H_0$}
\]

\section{ANOVA a una via}

Assumendo:
\begin{itemize}
    \item campioni indipendenti;
    \item normalità delle popolazioni;
    \item omoschedasticità (uguaglianza delle varianze);\newline
\end{itemize}

Si vuole testare:

\[
H_0:\mu_1=\mu_2=\cdots=\mu_k
\]

La statistica del test è:

$$
\scalebox{1.5}{$
\frac{
\frac{\sum_{i=1}^{k} n_i(\bar X_i-\bar X)^2}{k-1}
}{
\frac{\sum_{i=1}^{k}\sum_{j=1}^{n_i}(X_{ij}-\bar X_i)^2}{n-k}
}
$}
\overset{H_0}{\sim}
F_{k-1,n-k}
$$
$$
(\text{con }n_i=\text{numerosità del gruppo }i, \ \bar X_i\text{ media del gruppo }i, \ \bar X\text{ media generale})
$$

\[
\text{La regione critica è:}
\left\{
\text{Valore statistica}>
F_{1-\alpha,k-1,n-k}
\right\}
\]\newline

Esempio (5\%):

\begin{table}[ht]
    \centering
    \begin{tabular}{c|cccc}
     & $G_1$ & $G_2$ & $G_3$ & $G_4$ \\
    \hline
    1&279&378&172&381\\
    2&338&275&335&346\\
    3&334&412&335&340\\
    4&198&265&282&471\\
    5&303&286&250&318
    \end{tabular}
\end{table}

\[
\begin{aligned}
\bar X_1&=\frac{279+338+334+198+303}{5}=290.4\\
\bar X_2&=\frac{378+275+412+265+286}{5}=323.2\\
\bar X_3&=\frac{172+335+335+282+250}{5}=274.8\\
\bar X_4&=\frac{381+346+340+471+318}{5}=371.2\\
\bar X&=
\frac{1}{20}
\sum_{i=1}^{4}\sum_{j=1}^{5}X_{ij}
=314.9
\end{aligned}
\]

\[
\begin{aligned}
SS_B=
\sum_{i=1}^{4}n_i(\bar X_i-\bar X)^2&=
5[(290.4-314.9)^2+(323.2-314.9)^2\\ &+(274.8-314.9)^2+(371.2-314.9)^2]
\approx27234.20\\
SS_W
=
\sum_{i=1}^{4}
\sum_{j=1}^{5}
(X_{ij}-\bar X_i)^2&\approx63953.60
\end{aligned}
\]

\[
\frac{MS_B}{MS_W}
=\frac{SS_B/3}{SS_W/16}=
\frac{9078.07}{3997.10}
\approx2.27
\]

$$F_{0.95,3,16}\approx3.24>2.27\rightarrow\text{ Non rifiuto $H_0$}
$$

\section{Algoritmo EM}

\df L'algoritmo EM (Expectation-Maximization) è un metodo utilizzato per stimare i parametri di un modello quando sono presenti dati incompleti o variabili latenti.

Si vuole massimizzare la verosimiglianza dei dati osservati, ma la variabile latente rende difficile massimizzare direttamente, si considera quindi la verosimiglianza completa:

\[
L_c(\theta|X,Z)=p(X,Z|\theta)\ (X\text{ dati osservati},\
Z \text{ dati latenti},\
\theta\text{ parametri})
\]

\vspace{10pt}

L'algoritmo procede alternando due passi:

\[
\boxed{\text{E-step}\qquad Q(\theta|\theta^{(t)})
=
E_{Z|X,\theta^{(t)}}
[\log p(X,Z|\theta)]}
\]

\[
\boxed{\text{M-step}\qquad\theta^{(t+1)}
=
\arg\max_{\theta}
Q(\theta|\theta^{(t)})}
\]\newline

Fino a convergere:\footnote{La convergenza può essere verso un massimo locale, quindi la soluzione dipende dalla scelta iniziale $\theta^{(0)}$.}

\[
|\theta^{(t+1)}-\theta^{(t)}|<\varepsilon.
\]\newline

È garantito che ad ogni iterazione la verosimiglianza osservata non diminuisca:

\[
L(\theta^{(t+1)}|X)\geq L(\theta^{(t)}|X).
\]

\end{document}
